\documentclass[12pt,a4paper,onesided]{report}


%%%%%%%%%%%%%%%%%%%%%%%%%%%%%%%
% Foundation-Related Packages %
%%%%%%%%%%%%%%%%%%%%%%%%%%%%%%%
\usepackage[margin=2.5cm]{geometry}
%\usepackage{subfiles}

%%%%%%%%%%%%%%%%%%%%%%%%%%%%%%
%   Maths-Related Packages   %
%%%%%%%%%%%%%%%%%%%%%%%%%%%%%%
\usepackage[fleqn]{amsmath}
\usepackage{amsfonts}
\usepackage{amssymb}
\usepackage{siunitx}
\usepackage{fdsymbol}

%%%%%%%%%%%%%%%%%%%%%%%%%%%%%%%
%        SIUNITX setup        %
%%%%%%%%%%%%%%%%%%%%%%%%%%%%%%%
\sisetup{zero-decimal-to-integer, round-mode=places, round-precision=3}

%%%%%%%%%%%%%%%%%%%%%%%%%%%%%%%
%    Code-Related Packages    %
%%%%%%%%%%%%%%%%%%%%%%%%%%%%%%%
\usepackage{listings}
\usepackage[lined,boxed,procnumbered,longend,linesnumbered,commentsnumbered]{algorithm2e}

%%%%%%%%%%%%%%%%%%%%%%%%%%%%%%
%  Graphic-Related Packages  %
%%%%%%%%%%%%%%%%%%%%%%%%%%%%%%
\usepackage{graphicx}
\usepackage{pgfplots}
\usepackage{tikz}
\usepackage{xcolor}
\usepackage{xcolor-solarized}

%%%%%%%%%%%%%%%%%%%%%%%%%%%%%%
%       Other Packages       %
%%%%%%%%%%%%%%%%%%%%%%%%%%%%%%
\usepackage[style=ieee,backend=biber]{biblatex}
\usepackage{booktabs}
\renewcommand{\arraystretch}{1.3}

\usepackage{enumitem}
\usepackage{eurosym}
\usepackage{hyperref}
\usepackage{multirow}
\usepackage{parskip}
\usepackage{tocloft}
\usepackage{titlesec}
\usepackage{xparse}
\usepackage{enumitem}
\setlist[description]{leftmargin=10ex,labelindent=10ex}

%%%%%%%%%%%%%%%%%%%%%%%%%%%%%%
%  HEADER AND FOOTER CONFIG  %
%%%%%%%%%%%%%%%%%%%%%%%%%%%%%%
\usepackage{fancyhdr}
\usepackage{lastpage}

\pagestyle{fancy}

\lhead{}%\slshape \leftmark}
\rhead{\slshape \leftmark}
\rfoot{Page \thepage\ of \pageref{LastPage}}
\cfoot{}
%%%%%%%%%%%%%%%%%%%%%%%%%%%%%%%
%     SECTIONING CONFIG       %
%%%%%%%%%%%%%%%%%%%%%%%%%%%%%%%
%Chapter
\renewcommand{\chaptername}{\vspace{-3em} Task}
%\renewcommand{\thechapter}{\arabic{chapter}}
%\titleformat{\chapter}[display]   
%{\normalfont\huge\bfseries}{\chaptertitlename\ \thechapter}{20pt}{\Huge}   
%\titlespacing*{\chapter}{0pt}{-50pt}{40pt}


% Section
%\renewcommand{\thesection}{Task \arabic{chapter}.\arabic{section}} 
	
%Subsection
%\renewcommand{\thesubsection}{Task \arabic{chapter}\ifnum\value{section}>0.\arabic{section}\fi.\arabic{subsection}} 


%%%%%%%%%%%%%%%%%%%%%%%%%%%%%%%
%         TOC CONFIG          %
%%%%%%%%%%%%%%%%%%%%%%%%%%%%%%%
%\cftsetindents{chapter}{0em}{3.75em}
%\cftsetindents{section}{1.5em}{3em}
%\cftsetindents{subsection}{3em}{4.5em}

%\renewcommand{\cftchappresnum}{Task }

% Name to use for referenceing chapters
%\def\chapterautorefname{Task}

% listings configuration
\definecolor{mygreen}{rgb}{0,0.6,0}
\definecolor{mygray}{rgb}{0.5,0.5,0.5}
\definecolor{mymauve}{rgb}{0.58,0,0.82}

\lstset{inputpath="/home/mark/IdeaProjects/CSA1017/",% Folder storing code
}

\lstdefinestyle{prog}{
  backgroundcolor=\color{white},    % choose the background color; you must add \usepackage{color} or \usepackage{xcolor}
  basicstyle=\footnotesize\ttfamily\color{black},% the size of the fonts that are used for the code
  breakatwhitespace=false,          % sets if automatic breaks should only happen at whitespace
  breaklines=true,                  % sets automatic line breaking
  captionpos=b,                     % sets the caption-position to bottom
  commentstyle=\color{mygreen},     % comment style
  deletekeywords={...},             % if you want to delete keywords from the given language
  escapeinside={\%*}{*)},           % if you want to add LaTeX within your code
  extendedchars=true,               % lets you use non-ASCII characters; for 8-bits encodings only, does not work with UTF-8
  keepspaces=true,                  % keeps spaces in text, useful for keeping indentation of code (possibly needs columns=flexible)
  keywordstyle=\color{blue},        % keyword style
  language=Java,                    % the language of the code
  morekeywords={*,...},             % if you want to add more keywords to the set
  numbers=left,                     % where to put the line-numbers; possible values are (none, left, right)
  numbersep=5pt,                    % how far the line-numbers are from the code
  numberstyle=\tiny\color{mygray},  % the style that is used for the line-numbers
  rulecolor=\color{black},          % if not set, the frame-color may be changed on line-breaks within not-black text (e.g. comments (green here))
  showspaces=false,                 % show spaces everywhere adding particular underscores; it overrides 'showstringspaces'
  showstringspaces=false,           % underline spaces within strings only
  showtabs=false,                   % show tabs within strings adding particular underscores
  stepnumber=1,                     % the step between two line-numbers. If it's 1, each line will be numbered
  stringstyle=\color{mymauve},      % string literal style
  tabsize=2,                        % sets default tabsize to 2 spaces
}



\lstdefinestyle{sterm}{
    backgroundcolor=\color{solarized-base3},
    basicstyle=\ttfamily\small\color{solarized-base01},             % -- Making the font monospaced and small
    breaklines=true,                                                % -- Allows automatic line breaking
%   captionpos=b,                                                   % -- Caption position at bottom of listing
    commentstyle = \color{solarized-base1},                         % -- Sets comments to solarized-green
    escapeinside = {\%*}{*},                                        % -- Useful for adding LaTeX within the code
    emphstyle =\bfseries\color{solarized-cyan},                     % -- Emphasized text style to red
    identifierstyle=\color{solarized-base00},                       % -- Style for Identifiers
    keepspaces=true,                                                % -- Keeps Spaces in text, useful for keeping identation of code.
    keywordstyle =[1]\color{solarized-orange}\bfseries,             % -- Keyword style
    keywordstyle =[2]\color{solarized-yellow}\bfseries,             % -- Keyword2 style
    numbers=none,                                                   % -- Line numbers on the left
    showspaces=false,                                               % -- Do not show spaces everywhere by adding particular underscores
    showtabs=fase,                                                  % -- Do not show tabs within strings adding particular underscores
    stringstyle=\color{solarized-blue},                             % -- String style
    tabsize=2,                                                      % -- Setting the tab size to 2 spaces.
    framexleftmargin=0pt,
    framexrightmargin=0pt,
    framextopmargin=0pt,
}
\lstset{style=prog}

\tikzstyle{overbrace text style}=[font=\footnotesize, above, pos=.5, yshift=5mm]
\tikzstyle{overbrace style}=[decorate,decoration={brace,raise=2mm,amplitude=7.5pt}]
\tikzstyle{underbrace style}=[decorate,decoration={brace,raise=2mm,amplitude=3pt,mirror},color=gray]
\tikzstyle{underbrace text style}=[font=\tiny, below, pos=.5, yshift=-3mm]

\usetikzlibrary{decorations.pathreplacing}
\usetikzlibrary{arrows}
%%%%%%%%%%%%%%%%%%%%%%
%     NEW COMMANDS   %
%%%%%%%%%%%%%%%%%%%%%%
\newcommand{\source}{\pagebreak
	\ifnum\value{section}>0
		\subsection{The source code for \thesection}

	\else
		\section{The source code for Task \thechapter}

	\fi
	}
\newcommand{\AND}{{\normalfont\textbf{and}}}

\DeclareDocumentEnvironment{assumption}{o}{
  \begin{quotation}
  \centering
  \begin{tabular}{rp{0.65\linewidth}}
    Assumption\IfNoValueF{#1}{s} Made: &\em
  }{
  \end{tabular}
  \end{quotation}
  }

\DeclareDocumentEnvironment{limitation}{o}{
  \begin{quotation}
  \centering
  \begin{tabular}{rp{0.65\linewidth}}
    Limitation\IfNoValueF{#1}{s}: &\em
  }{
  \end{tabular}
  \end{quotation}
  }


\addbibresource{/home/mark/Documents/BibTeX/library.bib}

\begin{document}
	\begin{titlepage}
		\vspace*{\fill}
\centering
\includegraphics[width=0.2\textwidth]{University_of_Malta_Logo}\\
\vspace*{3em}
\textbf{\LARGE CSA 1017\\\vspace{0.1em}Data Structures and Algorithms 1}\\
\vspace*{0.75em}
\textbf{\Large Assignment}\\
\vspace{3em}
\large Mark Said Camilleri\\
B.Sc. (Hons) (Computing Science)\\
University of Malta\\
\vspace*{1em}
Friday 4th May 2016
\vfill
\end{titlepage}
\pagebreak
\tableofcontents
\addtocontents{toc}{\hfill\textbf{Page}\par}

\chapter*{Statement of Completion}
\addcontentsline{toc}{chapter}{Statement of Completion}
The questions below were the ones that have been attempted:
	\begin{description}
		\item[Question 1]This question has been successfully completed.
		\item[Question 2]This question has been successfully completed.
		\item[Question 3]This question has been successfully completed.
		\item[Question 4]This question has been successfully completed.
		\item[Question 5]This question has been successfully completed.
		\item[Question 6]This question has been successfully completed.
		\item[Question 7]This question has been successfully completed.
		\item[Question 8]This question has been successfully completed.
		\item[Question 9]This question has been successfully completed.
	\end{description}

\vspace{5em}
\begin{table}[h]
	\centering
\begin{tabular}{lp{20ex}p{20ex}lp{20ex}}
	 & & &  & \\\cline{1-2} \cline{4-5}
	 Signature & & & Date &\\
\end{tabular}
\end{table}
\chapter{Arabic to Roman Numeral Converter}
\vspace{-3em}
	\begin{tabular}{rp{36ex}p{36ex}}
		\toprule Input & Expected Output & Actual Output\\
		\midrule
		1 & I & I\\
		2 & II & II\\
		4 & IV & IV\\
		5 & V & V\\
		6 & VI & IV \\
		9 & IX & IX \\
		10 & X & X \\
		20 & XX & XX \\
		40 & XL & XL\\
		49 & XLIX & XLIX\\
		50 & L & L \\
		60 & LX & LX \\
		90 & XC & XC \\
		99 & XCIX & XCIX \\
		100 & C & C\\ 
		150 & CL & CL \\
		400 & CD & CD \\
		499 & CDXCIX & CDXCIX \\
		500 & D & D\\
		600 & DC & DC\\
		900 & CM & CM\\
		999 & CMXCIX & CMXCIX\\
		1000& M & M\\
		1024& MXXIV & MXXIV\\
		2000& The input was not a valid number between 1 and 1024 Please try again and enter a number to convert: &  The input was not a valid number between 1 and 1024 Please try again and enter a number to convert:\\
		abc& The input was not a valid number between 1 and 1024 Please try again and enter a number to convert: &  The input was not a valid number between 1 and 1024 Please try again and enter a number to convert:\\
				
		\bottomrule
				
	\end{tabular}
\source
\lstinputlisting{"/CSA1017 Q1/src/Question1.java"}
\chapter{Reverse Polish Notation Evaluator}
\begin{table}[h!]
	\centering
\begin{tabular}{lp{28ex}p{28ex}}
	\toprule Input & Expected Output & Actual Output\\
	\midrule
	4 5 + & 9.0 & 9.0\\
	12 6 - & 6.0 & 6.0\\
	3 2 / & 1.5 & 1.5 \\
	7 2 * & 14.0 & 14.0\\
	23 85 + 92 * & 9936.0 & 9936.0 \\
	43.5 3.2 - 4.5 * 3.24 + & 184.59 & 184.59 \\
	34.8 62.11 * -76 / & -28.43984211 & -28.439842105263157\\
	282 -56 * 1024 - -55.6 * & 934969.6 & 934969.6\\
	Test & Your expression contained invalid characters. For input string "T" Your expression is invalid. Evaluation failed& Your expression contained invalid characters. For input string "T" Your expression is invalid. Evaluation failed\\ 
	1 + & Stack is Empty. Your expression is invalid. Evaluation failed. &  Stack is Empty. Your expression is invalid. Evaluation failed.\\
	3 64 6 + & The stack has not been emptied. There are too many operands in your expression. Your expression is invalid. Evaluation failed. & The stack has not been emptied. There are too many operands in your expression. Your expression is invalid. Evaluation failed. \\
	4 0 / & Infinity & Infinity \\
	\bottomrule
\end{tabular}
\end{table}
\source
\subsection{Stack Class}
\lstinputlisting{"/CSA1017 Q2/src/Stack.java"}
\subsection{Question 2 Main Class}
\lstinputlisting{"/CSA1017 Q2/src/Question2.java"}
\chapter{Prime Numbers}
\section{Prime Number Checker using divisibility check}
\begin{table}[h!]
	\centering
\begin{tabular}{lp{30ex}p{30ex}}
	\toprule Input & Expected Output & Actual Output\\
	\midrule
	1 & 1 is not a prime number& 1 is not a prime number\\
	2 & 2 is a prime number & 2 is a prime number\\
	3 & 3 is a prime number & 3 is a prime number\\
	4 & 4 is not a prime number & 4 is not a prime number\\
	-2 & -2 is not a prime number & -2 is not a prime number\\a
	677 & 677 is a prime number & 677 is a prime number\\
	34939 & 34939 is a prime number & 34939 is a prime number\\
	188737204 & 188737204 is not a prime number & 188737204 is not a prime number\\
	381165334 & 381165334 is not a prime number & 381165334 is not a prime number\\
	947396057 & 947396057 is a prime number & 947396057 is a prime number\\
	9223372036854775807 & 9223372036854775807 is not a prime number & 9223372036854775807 is not a prime number\\
	-9223372036854775808 & -9223372036854775808 is not a prime number &-9223372036854775808 is not a prime number\\
	9223372036854775808 & Your input was not accepted. Please restart the program and try again with a valid input. & Your input was not accepted. Please restart the program and try again with a valid input.\\
	Test & Your input was not accepted. Please restart the program and try again with a valid input. & Your input was not accepted. Please restart the program and try again with a valid input.\\
	
	\bottomrule
	\end{tabular}
\end{table}
\source
\lstinputlisting{"/CSA1017 Q3/src/Question3a.java"}
\newpage
\section{Prime Number Checker using The Sieve of Eratosthenes}
For this an \lstinline|int| variable had to be used, due to it being the length of an array. As a result the range of numbers that can be checked was smaller.
\begin{table}[h!]
	\centering
	\begin{tabular}{lp{30ex}p{30ex}}
		\toprule Input & Expected Output & Actual Output\\
		\midrule
		1 & 1 is not a prime number& 1 is not a prime number\\
		2 & 2 is a prime number & 2 is a prime number\\
		3 & 3 is a prime number & 3 is a prime number\\
		4 & 4 is not a prime number & 4 is not a prime number\\
		-2 & Sieve of Eratosthenes only works on +ve integers. Please restart the program and try again with a valid input & Sieve of Eratosthenes only works on +ve integers. Please restart the program and try again with a valid input\\
		677 & 677 is a prime number & 677 is a prime number\\
		34939 & 34939 is a prime number & 34939 is a prime number\\
		188737204 & 188737204 is not a prime number & 188737204 is not a prime number\\
		381165334 & 381165334 is not a prime number & 381165334 is not a prime number\\
		947396057 & 947396057 is a prime number & 947396057 is a prime number\\
		2147483646& 2147483646 is not a prime number & Requested array size exceeds VM limit Please restart the program and try again with a smaller input\footnote{} \\ 
		-2147483646 & Sieve of Eratosthenes only works on +ve integers. Please restart the program and try again with a valid input & Sieve of Eratosthenes only works on +ve integers. Please restart the program and try again with a valid input\\
		9223372036854775807 & Your input was not accepted. Please restart the program and try again with a valid input.
		 & Your input was not accepted. Please restart the program and try again with a valid input.\\
		Test & Your input was not accepted. Please restart the program and try again with a valid input.
		& Your input was not accepted. Please restart the program and try again with a valid input.\\
		\bottomrule
	\end{tabular}
\end{table}
\footnotetext{This error depends on the amount of memory allocated for the program}
\source
\lstinputlisting{"/CSA1017 Q3/src/Question3b.java"}
\chapter{Shell Sort} 
Since the requirement of this task was for the program to generate an array containing 16,384 elements, all of which are random numbers, and sort them, it is impossible to write the contents of the array for a reliable amount of tests and keep this document to a reasonable size. As a result, the program itself checks whether the array that has been sorted, is in fact sorted. This is done using the algorithm \lstinline|checkSortedAscending| as given in Algorithm \ref{asc:alg} (also included in the source code).

\begin{algorithm}
	\TitleOfAlgo{Check if sorted in ascending order}
	\KwIn{int array}
	\KwOut{boolean}
	\BlankLine
		\For{i = 0... length of array - 1}{
			\If{element at i $>$ element at i+1}{
				\Return{false}
			}
		}	
		\Return{true}
	\caption{checkSortedAscending(int[] array)}
	\label{asc:alg}
\end{algorithm}

\source
\lstinputlisting{"/CSA1017 Q4/src/Question4.java"}
\chapter{}
\chapter{}
\chapter{}
\chapter{}
\chapter{}
\printbibliography[title={References}]
\addcontentsline{toc}{chapter}{References}
\end{document}
